\documentclass[journal]{IEEEtran}

\usepackage{amsmath}
\usepackage{graphicx}
\usepackage{cite}
\usepackage{url}

\hyphenation{op-tical net-works semi-conduc-tor}

\begin{document}

\title{Calibrated Structured Decisions for LLM-Based Applicant Assessment: An Empirical Study of System One Models in Candidate Ranking}

\author{Jessie~Angelica$^{1}$, Hasanul~Fahmi~Zuhri$^{1,2,*}$\\[0.5em]
{\normalsize $^{1}$Faculty of Computer Science, President University, Bekasi, Indonesia\\
$^{2}$Faculty of Artificial Intelligence \& Frontier Technologies, UNITAR International University, Petaling Jaya, Malaysia}%
\thanks{*Corresponding author: fahmi.zuhri@unitar.my.}}

\markboth{Calibrated Structured Decisions for LLM-Based Applicant Assessment}%
{Angelica \MakeLowercase{\textit{et al.}}: Calibrated Structured Decisions for LLM-Based Applicant Assessment}

\maketitle

\begin{abstract}
Large language model (LLM) pipelines for applicant ranking usually combine free-text assessment with iterative listwise comparison. This approach needs extra machinery to catch hallucinated claims, and it gives no built-in measure of how confident the model actually is. In September 2026, TypeSafe AI introduced Jev, the first publicly available ``System One Model.'' Jev is non-autoregressive: it answers a fixed set of typed questions in one parallel pass and reports a calibrated confidence with every answer, so it removes free-text generation -- and the hallucination risk that comes with it -- by design. No independent academic evaluation of this model class exists yet. The objective of this study is to close that gap: to propose and empirically validate three architecture-level adaptations that address concrete reliability problems observed when Jev is applied to applicant assessment without modification. The three adaptations are (i) criteria-grounded question design, which replaces bare ordinal labels with concrete situational descriptions to improve answer calibration; (ii) multi-call ensemble averaging, which aggregates independent single-pass calls to recover ranking stability that single-sample noise would otherwise destroy; and (iii) structured score decomposition, which uses the model's per-requirement granularity as a diagnostic signal for upstream pipeline defects that a monolithic free-text judgment cannot reveal. The research method is to rebuild an existing LLM-and-tournament applicant-ranking pipeline around Jev-based structured decisions -- replacing its free-text assessment and iterative Plackett-Luce tournament stages -- and to evaluate the resulting pipeline on a corpus of 501 resumes across 10 job profiles, under the same skill-matching shortlist configuration as the prior system, using four evaluation procedures: test-retest reliability, internal coherence, ranking convergence, and a criteria-design ablation. \emph{[Results placeholder: to be finalized once the full-corpus evaluation and significance testing are complete. A preliminary study on a 338-pair shortlisted subset shows that criteria-grounded question design raises per-requirement answer confidence from 0.705 to 0.902 (Wilcoxon signed-rank, $p < 0.001$, large effect), and that three-call ensemble averaging recovers ranking convergence from Kendall-$\tau$ = 0.906 to 0.945, approaching the 0.957 reported by the closest prior tournament-based system.]}
\end{abstract}

\begin{IEEEkeywords}
Applicant Ranking, Calibration, Ensemble Aggregation, Human Resources, Large Language Models, Structured Decision Models, System One Models.
\end{IEEEkeywords}

\IEEEpeerreviewmaketitle

\section{Introduction}
\IEEEPARstart{R}{ecruitment} platforms increasingly use large language models (LLMs) to read resumes, summarize applicant fit, and compare candidates against a job description, in place of manual screening or keyword search \cite{ref1,ref2,ref3,ref4}. These systems are useful for the same reason they are risky: an LLM generates free text, and free text can assert a strength or weakness that the source resume never actually supports. Prior work responds to this hallucination risk with prompt-level constraints, retry-on-contradiction loops, and post-hoc auditing \cite{ref2}. It also combines free-text assessment with iterative listwise tournaments and Plackett-Luce aggregation, turning per-applicant judgments into a defensible ranking \cite{ref1,ref2}. Both kinds of mitigation exist because the underlying generation mechanism -- next-token sampling over an open vocabulary -- gives no structural guarantee against an ungrounded claim, and no native signal of how confident the model really is.

\subsection{Research Objectives}
This paper asks whether a different generation paradigm -- one that answers predefined, typed questions instead of generating open text -- can replace the free-text assessment stage of an applicant-ranking pipeline without reintroducing the reliability problems that paradigm is meant to solve. Where the new paradigm introduces reliability problems of its own, we propose and empirically validate architecture-level adaptations that address them. Concretely, this paper (1) rebuilds an existing LLM-and-tournament applicant-ranking pipeline's assessment and ranking stages around Jev, TypeSafe AI's first System One Model; (2) identifies, through controlled empirical study on the pipeline's own corpus, three specific failure modes that appear when Jev is used without adaptation; and (3) proposes and validates a targeted architectural response to each one.

\subsection{Gap Analysis}
Three gaps motivate this study. First, System One Models are new enough -- Jev's public release predates this manuscript by days -- that no independent, peer-reviewed evaluation of the paradigm exists; every technical claim about its reliability currently comes from the vendor itself. Second, the recent literature on LLM calibration and self-consistency \cite{ref5,ref6,ref7} targets autoregressive judges that generate a scalar score or a free-text verdict. None of it addresses a model whose primitive output types are a bounded probability (Noul), a categorical distribution (Choice), or an ordinal distribution (Score), each answered in a single parallel pass -- so neither its diagnostic tools nor its mitigation strategies (e.g., resampling temperature, chain-of-thought prompting) transfer directly. Third, the LLM-hiring-fairness literature \cite{ref8,ref9,ref10} audits demographic bias in a system's final decision, but treats the underlying judgment mechanism as a black box. It does not ask whether that mechanism's own internal granularity could diagnose defects elsewhere in the pipeline -- defects that shaped which candidates reached the judgment stage in the first place.

\subsection{Contributions}
This paper makes three architecture-level contributions, each grounded in a controlled empirical finding rather than a purely theoretical argument:

\begin{enumerate}
\item \textbf{Criteria-Grounded Question Design}: a methodology for writing both Score- and Noul-type questions with concrete, situational evidence descriptions instead of bare ordinal labels or circular restatements of the question. A paired ablation shows this raises answer confidence for both question types, and a separate internal-coherence check shows it improves group separation for Score-derived judgments, on held-out (job, applicant) pairs.
\item \textbf{Multi-Call Ensemble Averaging}: an aggregation scheme that issues several independent Jev calls per candidate and combines them -- arithmetic mean for score-valued answers, confidence-weighted majority vote for categorical answers. This recovers most of the ranking-order stability of an iterative tournament without needing an iterative ranking process.
\item \textbf{Structured Score Decomposition as a Diagnostic Signal}: using the model's per-requirement score granularity, which a monolithic free-text judgment cannot offer, to detect and localize a pre-existing upstream shortlisting defect that a prose-only system could only expose through manual inspection.
\end{enumerate}

Section II reviews related work; Section III details the proposed pipeline and its three architectural contributions.

\section{Literature Review}

\subsection{LLM-Based Applicant Ranking Pipelines}
The system this paper extends combines LLM-driven free-text assessment with an active listwise tournament. Subsets of shortlisted applicants are repeatedly compared by an LLM, and the resulting partial orderings are aggregated into global utilities via a Bayesian Plackett-Luce model under Monte Carlo knowledge-gradient subset selection, following the design of Yuksel et al.~\cite{ref1}. That prior system also introduced a retry-on-contradiction loop that re-checks every generated weakness against the applicant's own previously extracted skills, positional-token identifiers that stop a locally hosted model from mis-copying applicant labels during ranking, and pool-size-adaptive iteration scaling \cite{ref2}. Other recent systems apply LLMs to resume extraction and pointwise scoring inside retrieval-augmented or multi-agent architectures, without a comparative ranking step \cite{ref3,ref4}. None of this body of work departs from free-text or free-form JSON generation as the model's fundamental output mode, so all of it inherits some version of the hallucination-mitigation burden that motivates \cite{ref2}'s retry loop.

\subsection{Calibration and Structured Decision Models}
A separate and rapidly growing literature asks whether an LLM's stated or implied confidence reflects its true probability of being correct. Yaldiz et al.~\cite{ref5} show that reinforcement learning from verifiable rewards improves task accuracy in decision-making LLMs, but also makes them noticeably overconfident relative to supervised fine-tuning; they propose a calibration-aware reinforcement objective that trades a small amount of accuracy for much better-calibrated confidence. Wang et al.~\cite{ref6} propose a structured-output reliability framework built on semantic tree-edit-distance and repeated-measurement consistency scoring, and find that reliability varies sharply across models and prompts even when the output schema is held fixed. Both studies, like the broader calibration literature they represent, treat the underlying model as autoregressive: confidence is estimated from output-token probabilities, or from resampling a free-text or JSON generation process. Jev's Noul, Choice, and Score primitives instead report confidence as a first-class field of a single non-autoregressive inference call, trained explicitly for that purpose under RLCD rather than estimated after the fact. This is a different enough mechanism that findings from the autoregressive-calibration literature cannot be assumed to transfer without independent verification -- which this paper undertakes for the hiring domain specifically.

\subsection{Self-Consistency and Aggregation for LLM Judgments}
Haldar and Hockenmaier~\cite{ref7} document substantial self-inconsistency in LLM-as-a-judge frameworks: the same judge, asked to re-evaluate the same input, frequently returns a different verdict, which undermines any pairwise or ranking comparison built on a single judgment. The usual mitigation in that literature is aggregation -- majority voting or averaging over multiple independent judgments to suppress pointwise noise rather than remove it at the source. This paper adopts the same aggregation principle, but applies it to a model whose single call is already an order of magnitude cheaper and faster than the autoregressive judges that motivated the technique. That makes multi-call aggregation a comparatively inexpensive architectural option here, rather than a costly last resort.

\subsection{Fairness and Bias in LLM-Based Hiring}
A substantial and independent line of work audits demographic bias in LLM-based hiring tools. Webster~\cite{ref8} finds that fairness and competence are separable failure modes: some resume-screening platforms that pass a narrow bias audit fail a resume-role competence test outright, so the two must be audited jointly. Wen et al.~\cite{ref9} and Seshadri et al.~\cite{ref10} each document persistent racial and gender disparities in LLM-driven resume evaluation and candidate retrieval, including in stages upstream of any final scoring decision. This literature evaluates a hiring system's final output against a demographic ground truth. It does not ask whether the judgment model's own internal structure can detect non-demographic defects -- such as an upstream retrieval stage admitting candidates who lack a role's defining skill -- which is the diagnostic use this paper proposes for Jev's structured score decomposition. Fairness auditing and pipeline-defect diagnosis are complementary concerns, not competing ones, and this paper's scope is limited to the latter.

\subsection{Positioning of This Work}
No paper reviewed here evaluates a System One Model, and no paper in the calibration or self-consistency literature addresses a non-autoregressive, single-pass, typed-question architecture. To the authors' knowledge, this is the first independent empirical study of this model class. It is situated in a domain -- applicant ranking -- already served by a directly comparable free-text predecessor \cite{ref1,ref2}, which makes an architecture-level comparison possible, not just a benchmark-level one.

\section{Methodology}

\subsection{Pipeline Overview}
The proposed pipeline keeps the ingestion, LLM-based skill extraction, and embedding-based shortlisting stages of the prior system \cite{ref2} unchanged. Job profiles are loaded as plain text and resumes are parsed from PDF; an LLM extracts a normalized, resume-grounded skill list for each job profile and applicant; every extracted skill is embedded into a shared matrix; and, per job profile, an applicant is shortlisted only if at least a minimum number of the profile's distinct technical skills are matched above a cosine-similarity threshold. The pipeline departs from the prior system at the next two stages. The prior system generated a free-text assessment (strengths, weaknesses, additional skills) and then ranked shortlisted applicants through an iterative Monte Carlo knowledge-gradient tournament with Plackett-Luce aggregation. The proposed pipeline instead issues one or more Jev calls per shortlisted (job, applicant) pair, each returning a structured assessment directly, and ranks applicants for a job profile by sorting on the resulting overall fit score -- removing the iterative ranking stage entirely.

\subsection{Structured Assessment via Jev}
Each Jev call receives a single textual state -- the job title, job description, the applicant's previously extracted skills, and the applicant's resume text -- together with a fixed set of typed questions answered against that state in one parallel pass. Every (job, applicant) pair, regardless of job profile, is asked three questions: a Score question estimating overall fit on a five-level scale, a Choice question selecting a hiring recommendation among ``hire,'' ``maybe,'' and ``no,'' and a Noul question estimating the probability that the applicant meets the job profile's stated minimum qualifications. One additional Score question is generated per distinct technical skill in the job profile's extracted requirement list, asking how well the applicant's resume evidences that specific requirement. The job-description extraction stage also identifies, where stated, named certifications, a seniority/experience requirement, and an education requirement; each yields one further Noul question (one per certification, one for seniority, one for education), asked only when the job description actually states that requirement, so a job profile with no certification or seniority language incurs no extra calls. This breaks the prior system's single, opaque ``meets minimum qualifications'' judgment into pieces that can be diagnosed separately -- a candidate failing due to a missing certification is now distinguishable from one failing due to insufficient experience -- while the original holistic question is kept unchanged alongside them, mirroring the relationship between the overall fit score and the per-requirement scores. Every answer carries a model-reported confidence, and Score-type answers additionally carry a probability distribution over the five levels; the reported score is mapped linearly onto a 0--100 scale for downstream ranking and reporting.

\subsection{Novelty 1: Criteria-Grounded Question Design}
An initial implementation defined every Score question's five levels with bare numeric labels (``0,'' ``25,'' ``50,'' ``75,'' ``100''). This design later turned out to yield much lower answer confidence than a revised design using concrete, situational level descriptions -- for example, distinguishing a requirement ``mentioned only as a bare skill-list item'' from one ``used substantively in a real role or project with clear responsibility.'' To isolate the effect of this change from ordinary sampling noise, we compared the two designs in a paired ablation. For a sampled subset of (job, applicant) pairs already assessed under the concrete-criteria design, we issued a second call using the original bare-label design against the identical state, then compared answer confidence pairwise across every shared question. We separated the comparison into three parts: the one criteria-affected overall-fit question; the criteria-unaffected recommendation and minimum-qualification questions, used as a negative control since their criteria never changed; and the many per-requirement questions, the primary criteria-affected quantity. A Wilcoxon signed-rank test with a rank-biserial effect size was applied to each part.

The same principle generalizes beyond Score questions. The seniority and education Noul questions introduced in Section III-B were initially given circular criteria -- ``true: the CV supports that the candidate meets this requirement,'' ``false: the CV does not support it'' -- which restates the question instead of naming evidence. An internal-coherence check (mean overall fit score for applicants answered \texttt{true} versus \texttt{false} on each Noul question, per job profile) found this circular phrasing produced weak, and in two cases reversed, separation between the two groups. The weak separation was made worse in the smallest job profiles by extreme class imbalance, as few as one or two applicants on one side of the split. Rewriting the criteria to name concrete evidence -- explicit years of experience, employment dates, or job titles for seniority; an explicitly stated degree, field of study, or credential for education -- follows the same logic as the Score-question fix above and needed no new validation method, only its re-application to a different question type. We ran a second paired ablation on the 169 (job, applicant) pairs whose job profile states a seniority or an education requirement: we reused the concrete-criteria confidence already on file for each pair, and issued one additional fresh call against the identical state with the original circular criteria. As Table~\ref{tab_noul_ablation} shows, the rewrite raises mean answer confidence for the pooled seniority/education comparison from 0.828 to 0.839 (Wilcoxon signed-rank, $n=212$ paired observations, $p=0.0063$, rank-biserial $r=0.179$), reaching significance for education alone ($n=146$, $p=0.033$) but not for seniority alone ($n=66$, $p=0.134$). The direction matches the Score-question result, but the size of the effect is an order of magnitude smaller: Score confidence separated the two designs almost completely, while here $r=0.179$. A binary Noul question likely leaves less room for a criteria rewrite to move confidence than a five-level Score question does. We rejected one alternative: generating criteria per job description with an LLM. That would reintroduce a free-text generation step -- and its hallucination risk -- into a pipeline whose whole point is removing exactly that, and it would make scores for the same skill (e.g., ``Python'') incomparable across job profiles that happened to receive differently worded criteria.

\begin{table}[!ht]
\caption{Seniority/education Noul criteria ablation: concrete evidence-based criteria vs. the original circular criteria, on the same 169 (job, applicant) pairs.}
\label{tab_noul_ablation}
\centering
\footnotesize
\begin{tabular}{lccc}
\hline
Comparison & $n$ & Mean confidence (concrete / circular) & $p$-value\\
\hline
Seniority & 66 & 0.799 / 0.790 & 0.134\\
Education & 146 & 0.857 / 0.845 & 0.033\\
Pooled & 212 & 0.839 / 0.828 & 0.0063\\
\hline
\end{tabular}
\end{table}

\subsection{Novelty 2: Multi-Call Ensemble Averaging}
Because each Jev call is independent and comparatively cheap, the pipeline issues $n$ independent calls per (job, applicant) pair -- three by default, matching the stability-repeat count of the prior tournament-based system \cite{ref2} -- and aggregates them into a single assessment. Numeric fields (overall fit score, per-requirement scores, and per-question confidence) are combined by arithmetic mean across the $n$ calls. The categorical recommendation is combined by majority vote, with ties broken in favor of the tied label carrying the higher mean confidence across the calls that selected it. Binary answers (minimum-qualification, certification, seniority, and education judgments) are combined by simple majority, defaulting to the conservative (negative) outcome on an exact tie. To check whether this aggregation actually improves ranking-order stability, rather than just averaging away genuine disagreement, we measured ranking convergence two ways on the same corpus: first from single, unaggregated calls (three independent per-pair scores, each ranked separately, with pairwise Kendall-$\tau$ computed across the resulting three rankings per job profile), and second from three independent three-call averages (nine total calls per pair, grouped into three trials of three, each trial's ranking built from its own three-call average, with pairwise Kendall-$\tau$ again computed across the three trial rankings).

\subsection{Novelty 3: Structured Score Decomposition as a Diagnostic Signal}
Because Jev reports a separate score for every extracted requirement rather than a single free-text verdict, a job profile's shortlisted applicant pool can be inspected at the level of individual skills rather than only at the level of an aggregate score. This paper uses that granularity diagnostically. For each job profile, we compute the signal-to-noise ratio of its shortlisted applicant pool: the standard deviation of applicants' mean overall-fit scores (between-applicant signal) divided by the mean within-applicant standard deviation across repeated calls (measurement noise). We correlate this ratio (Spearman) against that profile's ranking-convergence Kendall-$\tau$, to test whether ranking instability traces to genuine score-noise proximity between similarly qualified applicants, rather than to model unreliability. Separately, for job profiles with unusually low overall scores, we inspect the lowest-scoring shortlisted applicants' full per-requirement breakdowns directly, to determine whether low scores reflect genuinely weak applicants or a shortlisting-stage admission error -- for example, an applicant shortlisted for a role through generic, widely shared skills, while scoring zero on every skill that specifically defines that role.

\subsection{Experimental Setup}
The pipeline is evaluated on the same corpus and shortlist configuration as the prior system \cite{ref2}: 501 resumes across 10 synthetically authored job profiles, shortlisted under a cosine-similarity threshold of 0.8 and a minimum of five matched distinct technical skills per job profile. We apply four evaluation procedures to the resulting shortlisted (job, applicant) pairs: test-retest reliability, issuing three independent single calls per pair and measuring score variance and recommendation agreement across them; internal coherence, testing whether the model's own structured outputs agree with one another (per-requirement scores against the overall fit score via Spearman correlation; the categorical recommendation and the binary minimum-qualification answer against the overall fit score via Mann-Whitney U); ranking convergence, the Kendall-$\tau$ procedure described in Section III-D, computed both from single calls and from three-call averages; and the criteria-design ablation described in Section III-C. All statistical tests are two-sided, except where a directional hypothesis is stated explicitly (for example, that a ``hire'' recommendation carries a higher score than a ``no'' recommendation). Every reported effect size accompanies, rather than substitutes for, its corresponding significance test. As Table~\ref{tab_config} shows, these settings were held fixed across every evaluation run reported in this paper; only the shortlisted-pair count varies slightly run-to-run, because the skill-extraction and shortlisting stages are themselves LLM-driven and re-run independently for each corpus pass.

\begin{table}[!ht]
\caption{Pipeline configuration used for all runs reported in this paper.}
\label{tab_config}
\centering
\footnotesize
\begin{tabular}{p{1.5in}p{1.6in}}
\hline
Parameter & Value\\
\hline
Assessment model & Jev (\texttt{jev-latest}), TypeSafe AI\\
Skill/JD extraction model & Qwen2.5-14B-Instruct (4-bit, Ollama)\\
Skill embedding model & sentence-transformers/all-MiniLM-L6-v2\\
Corpus & 501 resumes, 10 job profiles\\
Shortlisted (job, applicant) pairs & 338--346 across 10 profiles, run-to-run\\
Skill-match threshold / min. matches & cosine $\geq$ 0.8 / $\geq$ 5 job skills\\
Independent calls averaged per assessment ($n$) & 3\\
Score question levels & 5 (concrete situational criteria)\\
Test-retest repeats per pair & 3\\
Ranking-convergence trials (averaged-score check) & 3 $\times$ 3 calls\\
Criteria-design ablation sample & 40 pairs\\
Ablation significance test & Wilcoxon signed-rank\\
\hline
\end{tabular}
\end{table}

\section{Conclusion}
This paper replaced the free-text assessment and iterative tournament-ranking stages of an existing LLM-based applicant-ranking pipeline with Jev, TypeSafe AI's first System One Model. We proposed three architecture-level adaptations -- criteria-grounded question design, multi-call ensemble averaging, and structured score decomposition as a diagnostic signal -- each motivated by a concrete reliability problem observed in practice, not by a purely theoretical concern. Criteria-grounded design, first validated for Score-type questions, also holds for Noul-type questions, though the benefit is much smaller there; Section III-C discusses why. Structured score decomposition proved useful beyond assessing applicants: because the model reports a score per individual requirement rather than one opaque judgment, its output can also expose upstream defects in the pipeline that produced the candidate pool, an audit path a free-text predecessor cannot offer without manual reading.

\subsection{Limitations}
Several limitations bound the claims made here. First, and most importantly, this study -- like the prior system it extends \cite{ref2} -- has not been validated against independent human-expert judgment. Every metric reported here measures internal reliability, coherence, or convergence, not hiring-decision correctness, and a highly self-consistent model can still be consistently wrong. Second, we performed no fairness or demographic-bias audit on the resulting rankings; the fairness literature reviewed in Section II-D addresses a different, complementary question that this paper does not attempt to answer. Third, measured end-to-end latency for realistic multi-question payloads (15--30 questions per call, driven by each job profile's extracted requirement count) was several times higher than the vendor's own published figures, which were obtained on payloads of three questions. This gap is consistent with per-call latency scaling with question count, but we did not isolate that relationship experimentally. Fourth, the internal-coherence, group-separation check for certification, seniority, and education questions (Section III-B) is weakened by small and imbalanced samples within some job profiles -- as few as one or two applicants on one side of a comparison. The paired-ablation confidence result for the seniority/education criteria rewrite does not share this weakness, since it compares the same 169 applicants against themselves under two criteria designs; but its effect size is small enough ($r=0.179$) that it should be read only as evidence the same design principle applies to both question types, not as proof the reliability gap between Noul and Score questions is closed.

\subsection{Future Work}
Three directions follow directly from the limitations above: an independent human-expert validation study, ideally using the same rubric-agreement methodology as the closest prior system \cite{ref1}; a demographic fairness audit of the resulting rankings, following the methodology established for free-text LLM hiring systems \cite{ref8,ref9,ref10} but adapted to a structured-decision model's output; and a targeted latency study isolating the relationship between question count per call and end-to-end response time, to determine whether question-batching or -splitting strategies can bring latency closer to the vendor's published figures without sacrificing the reliability gains reported here.

\begin{thebibliography}{10}

\bibitem{ref1}
K.~A. Yuksel, A.~B. Anees, A.~Elneima, S.~Hewavitharana, M.~Al-Badrashiny, and H.~Sawaf, ``Agentic AI for Human Resources: LLM-Driven Candidate Assessment,'' in \emph{Proc.\ 19th Conf.\ Eur.\ Chapter Assoc.\ Comput.\ Linguistics (EACL)}, vol.~3: System Demonstrations, Rabat, Morocco, Mar.~2026, pp.~341--348.

\bibitem{ref2}
J.~Angelica and H.~F.~Zuhri, ``Hallucination-Aware Self-Correction for LLM-Based Applicant Assessment,'' unpublished manuscript, 2026.

\bibitem{ref3}
A.~C. Tran and N.~C. Tran, ``A recruitment support system based on Large Language Models and Retrieval-Augmented Generation,'' \emph{Procedia Computer Science}, vol.~269, pp.~259--268, 2025.

\bibitem{ref4}
F.~P.-W. Lo, J.~Qiu, Z.~Wang, H.~Yu, Y.~Chen, G.~Zhang, and B.~Lo, ``AI Hiring with LLMs: A Context-Aware and Explainable Multi-Agent Framework for Resume Screening,'' arXiv:2504.02870, 2025.

\bibitem{ref5}
D.~N. Yaldiz, E.~Spiliopoulou, Z.~Qi, S.~Varia, S.~Doss, and N.~Pappas, ``Balancing Classification and Calibration Performance in Decision-Making LLMs via Calibration Aware Reinforcement Learning,'' arXiv:2601.13284, 2026.

\bibitem{ref6}
G.~Wang, J.~Yu, X.~Zhang, D.~Jiang, Y.~Song, T.~Deb, X.~Liu, and P.~He, ``STED and Consistency Scoring: A Framework for Evaluating LLM Structured Output Reliability,'' arXiv:2512.23712, 2025.

\bibitem{ref7}
R.~Haldar and J.~Hockenmaier, ``Rating Roulette: Self-Inconsistency in LLM-As-A-Judge Frameworks,'' in \emph{Findings of the Assoc.\ Comput.\ Linguistics: EMNLP 2025}, 2025.

\bibitem{ref8}
K.~T. Webster, ``Fairness Is Not Enough: Auditing Competence and Intersectional Bias in AI-powered Resume Screening,'' arXiv:2507.11548, 2025.

\bibitem{ref9}
A.~Wen, T.~Patil, A.~Saxena, Y.~Fu, S.~O'Brien, and K.~Zhu, ``FAIRE: Assessing Racial and Gender Bias in AI-Driven Resume Evaluations,'' arXiv:2504.01420, 2025.

\bibitem{ref10}
P.~Seshadri, H.~Chen, S.~Singh, and S.~Goldfarb-Tarrant, ``Small Changes, Large Consequences: Analyzing the Allocational Fairness of LLMs in Hiring Contexts,'' arXiv:2501.04316, 2025.

\bibitem{ref11}
TypeSafe AI, ``Introducing System One Models \& Jev,'' Company Blog, Sep.~2026. [Online]. Available: https://typesafe.ai/blog/introducing-system-one-models-and-jev

\bibitem{ref12}
TypeSafe AI, ``Quickstart,'' TypeSafe Developer Documentation. [Online]. Available: https://docs.typesafe.ai/introduction/quickstart

\end{thebibliography}

\end{document}
